% Lösungen Einheit 3 von 4 – Einheit 3 (zusammenbau v0.1)
\einheitenkopf{Einheit 3 von 4 · Einheit 3}

\erg{16}{a) rechts ($3x$) \quad b) $3x = 12$, $x = 4$ \quad c) $4x = 20$, $x = 5$ \quad d) $5x + 7 = 2x + 22$, $x = 5$ \quad e) $3x + 6 = x + 14$, $x = 4$ \quad f) $x + 8 = 2x - 2$, $x = 10$ \quad g) $2x + 6 = 2x + 6$, also $6 = 6$ – jede Zahl ist Lösung \quad h) $x = 7$; Probe: links $5 \cdot 5 = 25$, rechts $3 \cdot 7 + 4 = 25$ (wA)}
\erg{17}{a) $0{,}5x = 2{,}5$, $x = 5$}
\erg{18}{a) Die 5 hat kein $x$: $6x + 5 = 23$, $6x = 18$, $x = 3$. \quad b) Minus vor der Klammer dreht beide Vorzeichen: $15 - x + 3 = 2x$, $18 = 3x$, $x = 6$. \quad c) $3 = 3$ ist immer wahr: Jede Zahl ist Lösung, nicht nur 0. \quad d) Nein: Links steht eine Zahl plus 4, rechts dieselbe Zahl plus 9 – die Seiten sind nie gleich.}
